
\documentclass{article}
\usepackage{colortbl}
\usepackage{makecell}
\usepackage{multirow}
\usepackage{supertabular}

\begin{document}

\newcounter{utterance}

\centering \large Interaction Transcript for game `cladder', experiment `full\_v1.5\_default', episode 4386 with selfrefine{-}Qwen3.5{-}27B.
\vspace{24pt}

{ \footnotesize  \setcounter{utterance}{1}
\setlength{\tabcolsep}{0pt}
\begin{supertabular}{c@{$\;$}|p{.15\linewidth}@{}p{.15\linewidth}p{.15\linewidth}p{.15\linewidth}p{.15\linewidth}p{.15\linewidth}}
    \# & \multicolumn{2}{c}{Player} && \multicolumn{2}{c}{Game Master} \\
    \hline

    \theutterance \stepcounter{utterance}  
    & & & \multicolumn{4}{p{0.6\linewidth}}{
        \cellcolor[rgb]{0.9,0.9,0.9}{
            \makecell[{{p{\linewidth}}}]{
                \texttt{\tiny{[P1$\langle$GM]}}
                \texttt{You are an expert in causal inference. The following question is not a typical commonsense query, but rather a meticulously designed question created by a professor specializing in causal inference, intended to assess the students' mastery of the course content.} \\
\texttt{Imagine a self{-}contained, hypothetical world with only the following conditions, and without any unmentioned factors or causal relationships: Gender has a direct effect on smoking and peanut allergy. Smoking has a direct effect on tar deposit. Tar deposit has a direct effect on peanut allergy. Gender is unobserved. For nonsmokers and with no tar deposit, the probability of being allergic to peanuts is 35\%. For nonsmokers and with high tar deposit, the probability of being allergic to peanuts is 63\%. For smokers and with no tar deposit, the probability of being allergic to peanuts is 25\%. For smokers and with high tar deposit, the probability of being allergic to peanuts is 49\%. For nonsmokers, the probability of high tar deposit is 22\%. For smokers, the probability of high tar deposit is 73\%. For smokers, would it be more likely to see being allergic to peanuts if the person had been a nonsmoker?} \\
\texttt{Start your answer with 'Yes' or 'No', followed by additional reasoning or evidence to support your explanation.} \\
            }
        }
    }
     \\ \\

    \theutterance \stepcounter{utterance}  
    & \multicolumn{4}{p{0.6\linewidth}}{
        \cellcolor[rgb]{0.561,0.737,0.561}{
            \makecell[{{p{\linewidth}}}]{
                \texttt{\tiny{[P1$\rangle$GM]}}
                \texttt{No} \\
\\ 
\texttt{To determine whether a smoker would be more likely to be allergic to peanuts if they had been a nonsmoker, we must compare the observed probability of allergy for smokers, \$P(Y=1 | S=1)\$, with the counterfactual probability of allergy for smokers had they not smoked, \$P(Y=1 | \textbackslash{}text\{do\}(S=0))\$. Since Gender (\$G\$) is an unobserved confounder affecting both Smoking (\$S\$) and Peanut Allergy (\$Y\$), we cannot simply compare \$P(Y=1 | S=1)\$ with \$P(Y=1 | S=0)\$. Instead, we must use the provided conditional probabilities to calculate the specific risk for the smoker population under the intervention of not smoking.} \\
\\ 
\texttt{**1. Calculate the Observed Probability of Allergy for Smokers (\$P(Y=1 | S=1)\$)**} \\
\texttt{We use the Law of Total Probability, conditioning on Tar Deposit (\$T\$):} \\
\texttt{*   \$P(T=1 | S=1) = 0.73\$} \\
\texttt{*   \$P(T=0 | S=1) = 1 {-} 0.73 = 0.27\$} \\
\texttt{*   \$P(Y=1 | S=1, T=1) = 0.49\$} \\
\texttt{*   \$P(Y=1 | S=1, T=0) = 0.25\$} \\
\\ 
\texttt{\$\$P(Y=1 | S=1) = (0.49 \textbackslash{}times 0.73) + (0.25 \textbackslash{}times 0.27)\$\$} \\
\texttt{\$\$P(Y=1 | S=1) = 0.3577 + 0.0675 = 0.4252\$\$} \\
\\ 
\texttt{**2. Calculate the Counterfactual Probability for Smokers if they were Nonsmokers (\$P(Y=1 | \textbackslash{}text\{do\}(S=0))\$)**} \\
\texttt{The intervention \$\textbackslash{}text\{do\}(S=0)\$ sets the smoking status to 0. In this hypothetical scenario, the distribution of Tar Deposit changes to that of nonsmokers, but the direct effect of Gender on Allergy remains. However, the problem provides the conditional probabilities of \$Y\$ given \$S\$ and \$T\$. Under the intervention \$\textbackslash{}text\{do\}(S=0)\$, the probability of allergy is determined by the probabilities associated with the nonsmoking state (\$S=0\$) and the resulting tar distribution for nonsmokers.} \\
\\ 
\texttt{*   \$P(T=1 | S=0) = 0.22\$} \\
\texttt{*   \$P(T=0 | S=0) = 1 {-} 0.22 = 0.78\$} \\
\texttt{*   \$P(Y=1 | S=0, T=1) = 0.63\$} \\
\texttt{*} \\
            }
        }
    }
    & & \\ \\

    \theutterance \stepcounter{utterance}  
    & & & \multicolumn{2}{p{0.3\linewidth}}{
        \cellcolor[rgb]{0.9,0.9,0.9}{
            \makecell[{{p{\linewidth}}}]{
                \texttt{\tiny{[GM$|$GM]}}
                \texttt{No} \\
            }
        }
    }
    & & \\ \\

    \theutterance \stepcounter{utterance}  
    & & & \multicolumn{2}{p{0.3\linewidth}}{
        \cellcolor[rgb]{0.9,0.9,0.9}{
            \makecell[{{p{\linewidth}}}]{
                \texttt{\tiny{[GM$|$GM]}}
                \texttt{no} \\
            }
        }
    }
    & & \\ \\

    \theutterance \stepcounter{utterance}  
    & & & \multicolumn{2}{p{0.3\linewidth}}{
        \cellcolor[rgb]{0.9,0.9,0.9}{
            \makecell[{{p{\linewidth}}}]{
                \texttt{\tiny{[GM$|$GM]}}
                \texttt{game\_result = WIN} \\
            }
        }
    }
    & & \\ \\

\end{supertabular}
}

\end{document}
